% ==========================================================
% STANDARDIZED PREAMBLE (STATIC CORE) — DO NOT EDIT THIS PART
% ==========================================================

\documentclass[11pt,a4paper]{article}

\usepackage{iftex}
\ifPDFTeX
  \usepackage[utf8]{inputenc}
  \usepackage[T1]{fontenc}
  \usepackage{stix2}
\else
  \usepackage{fontspec}
  \usepackage{unicode-math}
  \setmainfont{STIX Two Text}
  \setmathfont{STIX Two Math}
\fi

\usepackage[margin=2.5cm]{geometry}
\usepackage{parskip}
\usepackage{microtype}
\usepackage{enumitem}

\usepackage{amsmath, amssymb}

\usepackage{tocloft}
\setlength{\cftbeforesecskip}{2pt}
\setlength{\cftbeforesubsecskip}{1pt}

\widowpenalty=10000
\clubpenalty=10000

\usepackage{xcolor}
\definecolor{myblue}{RGB}{0, 51, 153}

\usepackage{sectsty}
\partfont{\color{myblue}}
\chapterfont{\color{myblue}}
\sectionfont{\color{myblue}}
\subsectionfont{\color{myblue}}
\subsubsectionfont{\color{myblue}}

\usepackage{hyperref}

\renewcommand{\cftsecfont}{\bfseries\color{myblue}}
\renewcommand{\cftsubsecfont}{\color{myblue}}
\renewcommand{\cftsubsubsecfont}{\color{myblue}}
\renewcommand{\cftsecpagefont}{\bfseries}
\renewcommand{\cftsubsecpagefont}{}
\renewcommand{\cftsubsubsecpagefont}{}

\hypersetup{
  colorlinks=true,
  linkcolor=myblue,
  citecolor=myblue,
  urlcolor=myblue
}

\usepackage{titlesec}
\titlespacing{\section}{0pt}{12pt}{6pt}
\titlespacing{\subsection}{0pt}{10pt}{5pt}
\titlespacing*{\section}{0pt}{2.5ex plus 1ex minus .2ex}{1.5ex plus .2ex}
\titlespacing*{\subsection}{0pt}{2ex plus 0.8ex minus .2ex}{1.2ex plus .2ex}
\titlespacing*{\subsubsection}{0pt}{1.8ex plus 0.6ex minus .2ex}{1ex plus .2ex}

\makeatletter
\newcommand{\StartBlueDoc}[3][]{%
  \title{\vspace{-2em}\textbf{#2}%
    \if\relax\detokenize{#3}\relax\else\\[2pt]{\large #3}\fi
    \vspace{+0.6em}}
  \author{#1}
  \date{\vspace{-0.6em}\today}
  \begin{document}
  \maketitle
  \tableofcontents
  \vspace{1em}\hrule\vspace{1em}
}
\makeatother

% ==========================================================
% CUSTOMIZATION ZONE
% ==========================================================

\newcommand{\DocAuthor}{}
\newcommand{\DocTitle}{Types of Research Software Tests}
\newcommand{\DocSubtitle}{}

\setcounter{tocdepth}{2}

\usepackage{booktabs}
\usepackage{tabularx}
\usepackage{xltabular}
\newcolumntype{L}{>{\raggedright\arraybackslash}X}
\newcolumntype{T}{>{\raggedright\arraybackslash}p{0.26\textwidth}}
\newcolumntype{P}[1]{>{\raggedright\arraybackslash}p{#1\textwidth}}

% Header rows of the table "At a glance": a section, with its number and its
% title.
\newcommand{\GlanceSection}[1]{%
  \multicolumn{3}{@{}l}{\textbf{\ref{#1}\quad\nameref{#1}}}}

% The card that describes a type. Every type has the same fields, in this
% order: \Definition, \Catches, \Misses, \Cost, \PairsWith, \Example, \See.
\newenvironment{typecard}
  {\begin{description}[style=multiline, leftmargin=6.5em,
     font=\normalfont\bfseries, topsep=2pt, itemsep=2pt,
     parsep=0pt]}
  {\end{description}}
\newcommand{\Definition}{\item[Definition]}
\newcommand{\Catches}{\item[Catches]}
\newcommand{\Misses}{\item[Misses]}
\newcommand{\Cost}{\item[Cost]}
\newcommand{\PairsWith}{\item[Pairs with]}
\newcommand{\Example}{\item[Example]}
\newcommand{\See}{\item[See]}

% Cite a principle of the other documents by its title. A \ref cannot
% reach into another PDF, so the title is the reference.
\newcommand{\DesignRef}[1]{\emph{Design Principles}, ``#1''}
\newcommand{\TestingRef}[1]{\emph{Testing Research Software}, ``#1''}

% ==================
% Start the document
% ==================
\StartBlueDoc[\DocAuthor]{\DocTitle}{\DocSubtitle}



% ==========================================================
\section*{About this document}
\phantomsection
\addcontentsline{toc}{section}{About this document}
% ==========================================================

This document describes the types of tests for research software: what each
type is, what it catches, and what it misses. It belongs to the same set as
\emph{Design Principles for Research Software}, cited below as \emph{Design
Principles}, and \emph{Testing Research Software}, and it is the one place
where the set defines the terms of testing. The other two documents use them
with the meaning given here.

Every test makes three choices: its scope, what runs
(Section~\ref{sec:scope}); its inputs, what goes in
(Section~\ref{sec:inputs}); and its oracle, what the result is checked
against (Section~\ref{sec:oracles}). Each subsection of these three sections
describes one type, and a test takes one type from each section. A
property-based test, for example, is usually a unit test on generated inputs
with an invariant as its oracle. The three choices are largely independent:
a unit test can use any oracle, and most oracles can check a single function
as well as a whole workflow. Some combinations do not work, though, such as
generated inputs with a known answer, because an arbitrary input has none;
the field ``Pairs with'' of each type says which.

A test checks one claim, and each oracle sees some errors and is blind to
others. One test, however well chosen, therefore checks only one aspect of
an entry point, and an entry point is tested by several tests with different
oracles. How many, and which, is the subject of \TestingRef{A test is only as
good as its oracle}.

Many tests are known by a name that is not a type of its own, such as smoke
test or regression test. Such a name stands for a combination of types, or
for the purpose a test serves; Section~\ref{sec:by-name} lists common ones.

This document states no rules. Where a type comes with a requirement or a
default, such as the tolerance of a comparison of floating-point numbers,
its field ``See'' cites the principle that states it.



\subsection*{How to read a type}
\phantomsection
\addcontentsline{toc}{subsection}{How to read a type}
\label{sec:fields}

Every type is described by the same fields, in the same order:

\begin{description}
  \item[Definition] What the type is.
  \item[Catches] The errors it finds.
  \item[Misses] The errors it cannot find, and so the reason to combine it
    with other types.
  \item[Cost] The effort to write it, the time to run it, and the work to
    keep it up to date.
  \item[Pairs with] The types of the other two sections that it works with,
    and those it does not.
  \item[Example] One test of the type, from a scientific field.
  \item[See] The principles that state rules about the type, and related
    types.
\end{description}



\subsection*{Terminology}
\phantomsection
\addcontentsline{toc}{subsection}{Terminology}
\label{sec:terminology}

The terms defined in \emph{Design Principles}, such as package, library code,
workflow code, entry point, host stack, and caller, carry the same meaning
here. The terms of testing are defined here, for all three documents:

\begin{description}
  \item[Test] Code that runs part of the package on an input and checks the
    result automatically, with the outcome pass or fail.
  \item[Test suite] All tests of a package, run together by the test runner
    of the host stack, such as testthat in R or pytest in Python.
  \item[Oracle] The source of the expected value that a test compares a
    result with, or of the property that it checks
    (Section~\ref{sec:oracles}).
  \item[Tolerance] The largest difference between a result and its expected
    value that a test accepts.
\end{description}

Each type, such as unit test, edge case, or snapshot, is defined in its own
subsection of Sections~\ref{sec:scope} to~\ref{sec:oracles}, and the common
names of tests, such as smoke test, in Section~\ref{sec:by-name}.



% The overview gets a page of its own, so that it can be printed.
\clearpage
% ==========================================================
\section*{At a glance}
\phantomsection
\addcontentsline{toc}{section}{At a glance}
% ==========================================================

Every type of this document in one line. The number and the title link to
the type. A test takes one type from each of the three sections.

{\small
\begin{tabularx}{\textwidth}{@{}lTL@{}}
  \toprule
  \GlanceSection{sec:scope} \\
  \midrule
  \ref{sec:unit} & \nameref{sec:unit}
    & One function, or a few that work closely together, called directly. \\
  \ref{sec:integration} & \nameref{sec:integration}
    & Several parts run together, to check that they fit. \\
  \ref{sec:end-to-end} & \nameref{sec:end-to-end}
    & An entry point run the way its caller runs it, from start to
      finish. \\
  \addlinespace
  \GlanceSection{sec:inputs} \\
  \midrule
  \ref{sec:examples} & \nameref{sec:examples}
    & Small, typical inputs that the author of the test writes into it. \\
  \ref{sec:edges} & \nameref{sec:edges}
    & Empty, single, tied, missing, extreme, and nearly singular inputs. \\
  \ref{sec:generated} & \nameref{sec:generated}
    & Many inputs, produced by a program from the documented input
      domain. \\
  \ref{sec:simulated} & \nameref{sec:simulated}
    & Data drawn from the model, with parameters known exactly. \\
  \ref{sec:real-data} & \nameref{sec:real-data}
    & Data from a real measurement or collection, or an excerpt of it. \\
  \addlinespace
  \GlanceSection{sec:oracles} \\
  \midrule
  \ref{sec:it-runs} & \nameref{sec:it-runs}
    & The code finishes without an error; the result is not checked. \\
  \ref{sec:snapshot} & \nameref{sec:snapshot}
    & The result agrees with a stored output of an earlier run. \\
  \ref{sec:reproducibility} & \nameref{sec:reproducibility}
    & The result agrees with that of another run of the same code. \\
  \ref{sec:known-answer} & \nameref{sec:known-answer}
    & The result agrees with an exact value obtained without the code. \\
  \ref{sec:invariant} & \nameref{sec:invariant}
    & The result has a property that every correct result has. \\
  \ref{sec:metamorphic} & \nameref{sec:metamorphic}
    & A known change of the input changes the result in a known way. \\
  \ref{sec:round-trip} & \nameref{sec:round-trip}
    & An operation followed by its inverse returns the input. \\
  \ref{sec:second-opinion} & \nameref{sec:second-opinion}
    & The result agrees with that of an independent implementation. \\
  \ref{sec:recovery} & \nameref{sec:recovery}
    & The method recovers the parameters behind simulated data, with honest
      uncertainty. \\
  \ref{sec:ground-truth} & \nameref{sec:ground-truth}
    & The result agrees with a truth about real data that was established
      without the method. \\
  \ref{sec:convergence} & \nameref{sec:convergence}
    & The error falls with the resolution at the rate of the method. \\
  \bottomrule
\end{tabularx}}



% ==========================================================
\section{Scope: what runs}
\label{sec:scope}
% ==========================================================

The scope of a test is the part of the package that it runs. The smaller the
scope, the faster the test and the more precisely a failure points to its
cause; the larger the scope, the more of the connections between the parts
the test checks.

\subsection{Unit test}
\label{sec:unit}

\begin{typecard}
  \Definition A test that calls one function, or a few functions that work
    closely together, directly and apart from the rest of the package.
  \Catches Errors in the logic of that function, such as a wrong formula,
    sign, index, or convention. A failing unit test points to the function
    that broke.
  \Misses Errors in how the parts fit together, such as a result passed on in
    the wrong units or order, a file that a later step cannot read, or an
    option that never reaches the function.
  \Cost The lowest of the three scopes: on a tiny input, a unit test
    finishes in a fraction of a second, so hundreds of them run in minutes.
    It needs code split into functions that can be called on their own.
  \PairsWith Every type of input and of oracle. Hand-picked examples
    (\ref{sec:examples}) with known answers (\ref{sec:known-answer}) make the
    classic unit test; generated inputs (\ref{sec:generated}) with
    invariants (\ref{sec:invariant}) make a property-based test.
  \Example The log-likelihood of three observations under a normal
    distribution, compared with the value worked out by hand.
  \See \TestingRef{Many small tests, few big ones}.
\end{typecard}



\subsection{Integration test}
\label{sec:integration}

\begin{typecard}
  \Definition A test that runs several parts together, to check that they
    fit: the result of one function passed to the next, a file written by one
    step and read by another, or a call into an external library, such as a
    solver or a reader of a file format. Its scope lies between that of a
    unit test and that of an end-to-end test.
  \Catches Mismatches at the interfaces, such as units or axis orders that
    differ between the parts, identifiers lost on the way, and assumptions
    about a dependency that do not hold, such as the meaning of its tolerance
    argument.
  \Misses Errors inside a part that the combination happens to hide, and
    errors in the configuration and the glue of a whole entry point.
  \Cost Higher than for a unit test: it runs longer, it needs inputs that
    suit every part, and a failure points to several parts at once.
  \PairsWith Hand-picked examples (\ref{sec:examples}) and simulated data
    (\ref{sec:simulated}); invariants (\ref{sec:invariant}), such as counts
    and identifiers that survive each interface; round trips
    (\ref{sec:round-trip}) through a file format.
  \Example The peaks found by a peak detector, passed to the function that
    fits them, keep the identifiers and the units of their detection.
  \See \DesignRef{Derive abstractions from applications, not applications
    from abstractions}, on keeping the driving applications as integration
    tests.
\end{typecard}



\subsection{End-to-end test}
\label{sec:end-to-end}

\begin{typecard}
  \Definition A test that runs an entry point the way its caller runs it,
    from start to finish: a workflow or a command, from its input files and
    configuration to its outputs, or, for library code, the calls of a
    complete analysis, such as a documented example.
  \Catches Steps that do not fit together, files that a step cannot find or
    read, options of the configuration that never reach their step, and
    errors in code that no unit test runs.
  \Misses Where an error is: a failure says that something broke, not what.
    Also every input and branch that the small dataset does not reach.
  \Cost Minutes rather than seconds, and a small dataset that reaches every
    step and every branch of the configuration.
  \PairsWith Simulated data (\ref{sec:simulated}) with parameter recovery
    (\ref{sec:recovery}) as its oracle; invariants (\ref{sec:invariant}),
    such as the number of samples and the total of the counts; snapshots
    (\ref{sec:snapshot}) of the outputs. With ``it runs''
    (\ref{sec:it-runs}) as its only oracle, it is a smoke test
    (Section~\ref{sec:by-name}).
  \Example A workflow run on six simulated samples: every sample appears in
    the final table, and the estimated fold changes recover the simulated
    ones.
  \See \DesignRef{Test the whole workflow on a small dataset};
    \TestingRef{Many small tests, few big ones}.
\end{typecard}



% ==========================================================
\section{Inputs: what goes in}
\label{sec:inputs}
% ==========================================================

The inputs decide which part of the input space a test explores. Inputs
that a person chooses are few and go where the person looks, which is rarely
where the bugs are. Inputs that a program generates reach much more of the
space, but an arbitrary input comes without a known answer, so the inputs
also decide which oracles a test can use.

\subsection{Hand-picked examples}
\label{sec:examples}

\begin{typecard}
  \Definition Small inputs that the author of a test chooses and writes into
    it, usually typical ones: a few observations, a short sequence, a
    three-by-three matrix.
  \Catches Errors on typical inputs, which are the inputs that most users
    bring. A small input makes the claim of the test readable and its
    expected value easy to derive.
  \Misses Everything the author did not think of. People choose inputs that
    they expect to work, so hand-picked examples rarely reach the inputs
    where code breaks.
  \Cost Low to write and to run.
  \PairsWith Known answers (\ref{sec:known-answer}) worked out by hand;
    every scope, and most other oracles.
  \Example A two-by-two table of counts, small enough that the p-value of
    Fisher's exact test can be computed by hand.
  \See Edge cases (\ref{sec:edges}) and generated inputs
    (\ref{sec:generated}), which reach what hand-picked examples miss.
\end{typecard}



\subsection{Edge cases}
\label{sec:edges}

\begin{typecard}
  \Definition Inputs at the edges of the input space: in size, such as an
    empty input and a single element; in value, such as zeros, ties, and the
    boundaries of a parameter's range; missing and special values, such as
    NA, NaN, and infinities; extreme magnitudes and nearly singular inputs;
    and unusual structure, such as unsorted input.
  \Catches Crashes, silent NaNs, off-by-one errors, divisions by zero, and
    wrong handling of missing values: code breaks at the edges, and typical
    data never go there. Writing the test also reveals edges where the
    behavior is not documented.
  \Misses Errors in the interior of the input space, on typical data.
  \Cost Low: the edges of most functions follow from a short list, and each
    test is fast.
  \PairsWith Known answers (\ref{sec:known-answer}), including the
    documented error or warning as the expected outcome; invariants
    (\ref{sec:invariant}). Generated inputs (\ref{sec:generated}) search the
    edges automatically.
  \Example The correlation of a variable with zero variance, which is
    undefined: the result is the documented one, such as NA, or the
    documented error.
  \See \TestingRef{Most bugs live at the edges}; \DesignRef{Raise, warn, or
    stay silent}; \DesignRef{Mind the gap, don't fill it}.
\end{typecard}



\subsection{Generated inputs}
\label{sec:generated}

\begin{typecard}
  \Definition Many inputs, produced by a program from a description of the
    input domain: random draws, systematic grids, or the inputs of a
    property-based testing framework, which also searches the edges and
    reduces a failing input to a minimal one.
  \Catches Failures on inputs that no one thought of, in rare regions of the
    input domain and in unusual combinations of values.
  \Misses Errors that the oracle cannot see, because an arbitrary input has
    no known answer, and inputs outside the domain that the generator
    describes.
  \Cost Moderate: describing the input domain and finding an oracle that
    holds for every input take thought, and a test runs the code hundreds of
    times.
  \PairsWith Invariants (\ref{sec:invariant}), metamorphic relations
    (\ref{sec:metamorphic}), round trips (\ref{sec:round-trip}), and second
    opinions (\ref{sec:second-opinion}), which need no known answer; ``it
    runs'' (\ref{sec:it-runs}), in a fuzz test. Known answers
    (\ref{sec:known-answer}) only where a formula gives the answer for every
    input.
  \Example Data frames of random size and content, generated by a
    property-based testing framework, on which reordering the rows leaves an
    estimate unchanged.
  \See \TestingRef{Let the computer pick the inputs}; \TestingRef{Flaky
    tests cry wolf}, on randomness in tests.
\end{typecard}



\subsection{Simulated data}
\label{sec:simulated}

\begin{typecard}
  \Definition Data drawn from the model that the method assumes, with
    parameters chosen by the test, so that the truth behind the data is
    known exactly. Unlike generated inputs, simulated data look like the
    data of an application, in size and in structure.
  \Catches Errors of the whole method on data of realistic size and
    structure, such as bias and standard errors that are too small.
  \Misses Failures on data that depart from the model, such as outliers,
    artifacts, and the quirks of real files, and every error that the
    simulator shares with the method, such as the same wrong
    parameterization.
  \Cost Moderate to high: the simulator is written from the definition of
    the model, apart from the code of the method, and studies over many
    simulated datasets are slow.
  \PairsWith Parameter recovery (\ref{sec:recovery}), which needs it;
    invariants (\ref{sec:invariant}); end-to-end tests
    (\ref{sec:end-to-end}) of workflows, whose small dataset it can provide
    with a known result.
  \Example Counts for a few hundred genes, drawn from a negative binomial
    distribution with a known mean and dispersion for each gene.
  \See \TestingRef{Plant the answer, then find it}.
\end{typecard}



\subsection{Real data}
\label{sec:real-data}

\begin{typecard}
  \Definition Data from a real measurement, experiment, or collection, or an
    excerpt of it small enough for the repository.
  \Catches What simulated data leave out: quirks of file formats, codes for
    missing values, unexpected identifiers, outliers, artifacts of
    instruments, and departures from the assumptions of the model.
  \Misses Whether the result is right, unless the data come with a known
    answer or a ground truth. Without one, a test on real data shows only
    that the code runs, or that its result has not changed. Where there is a
    truth, it is often uncertain and covers only part of the data.
  \Cost Finding data that are small, licensed for redistribution, and not
    confidential. Larger datasets make slow tests.
  \PairsWith Ground truth (\ref{sec:ground-truth}), where the data have one;
    known answers (\ref{sec:known-answer}), where a result of the
    computation is certified, as for the Longley data among the Statistical
    Reference Datasets of NIST; otherwise ``it runs'' (\ref{sec:it-runs}),
    snapshots (\ref{sec:snapshot}), invariants (\ref{sec:invariant}), and
    metamorphic relations (\ref{sec:metamorphic}). With a known signal
    planted into them, real data serve parameter recovery
    (\ref{sec:recovery}) as well.
  \Example A few hundred lines of a public file of weather station records,
    with its codes for missing values and its irregular timestamps, to test
    the reader of the format.
  \See \TestingRef{Every test stands on its own}, on licenses and
    confidential data.
\end{typecard}



% ==========================================================
\section{Oracle: what the result is checked against}
\label{sec:oracles}
% ==========================================================

The oracle decides what a passing test shows. The first three oracles check
only that the code runs, or compare the code with itself: they show that a
result exists or has not changed, not that it is right. The others check the
result against something independent of the code under test, and each sees a
different kind of error. \TestingRef{A test is only as good as its oracle}
lists strong oracles for common kinds of code.

\subsection{It runs}
\label{sec:it-runs}

\begin{typecard}
  \Definition The test passes if the code finishes without an error. Its
    result is not checked, or only its form, such as the number of output
    files or the shape of an array.
  \Catches Crashes, missing dependencies, broken installations, and steps
    that do not fit together because a file is missing or has the wrong
    format.
  \Misses Every wrong result that does not crash, which is the typical
    failure of research software: the plausible, wrong number.
  \Cost The lowest of all oracles.
  \PairsWith End-to-end tests (\ref{sec:end-to-end}) on a small hand-picked
    example, as a smoke test, and generated inputs (\ref{sec:generated}), as
    a fuzz test (Section~\ref{sec:by-name}). It adds little to a unit test,
    whose result is cheap to check.
  \Example A workflow run on a small dataset finishes and writes every output
    file it promises.
  \See \DesignRef{Test the whole workflow on a small dataset}, which asks for
    more than this oracle: it requires checking the outputs, not only that
    the run finished.
\end{typecard}



\subsection{Snapshot}
\label{sec:snapshot}

\begin{typecard}
  \Definition The result is compared with a reference output stored from an
    earlier run of the same code. A test with this oracle is a snapshot
    test.
  \Catches Silent changes, such as numbers that move after a refactoring, an
    update of a dependency, or a new compiler. For a large output, such as a
    table, a file, or a figure, it is often the only practical test.
  \Misses Whether the result is right. A snapshot shows only that the result
    has not changed, and a snapshot of a wrong result protects the error.
  \Cost Low to create. Every intended change of the result needs a decision
    and an update of the reference output, and a comparison bit for bit
    breaks on every new platform.
  \PairsWith Every scope; fixed inputs, such as hand-picked examples
    (\ref{sec:examples}) and real data (\ref{sec:real-data}). It shows more
    when the reference output was checked against another oracle before it
    was stored.
  \Example The table of model estimates for a reference dataset, compared
    within a tolerance with the table stored at the last release.
  \See \TestingRef{Unchanged is not correct}; \TestingRef{Compare floats
    with tolerance}.
\end{typecard}



\subsection{Reproducibility}
\label{sec:reproducibility}

\begin{typecard}
  \Definition The result is compared with that of another run of the same
    code on the same input: a run with the same seed, with another number of
    threads or processes, on another platform, or with other versions of the
    dependencies.
  \Catches Hidden nondeterminism, such as an unset seed, a race condition,
    uninitialized memory, or a dependence on global state or on the order in
    which files are listed; and results that change across platforms or
    versions of the dependencies by more than the tolerance.
  \Misses Every error that is the same in each run: a wrong result,
    reproduced, is still wrong.
  \Cost Low: the code runs twice. A comparison across platforms needs
    automatic test runs on each of them.
  \PairsWith Every scope and every type of input; it matters most for
    stochastic and parallel code.
  \Example A Monte Carlo simulation run twice with the same seed, then on one
    and on eight threads, with the results compared within the tolerance
    that defines reproducibility.
  \See \DesignRef{Once is luck, twice is science}; \TestingRef{Flaky tests
    cry wolf}.
\end{typecard}



\subsection{Known answer}
\label{sec:known-answer}

\begin{typecard}
  \Definition An exact result for a particular input, obtained without the
    code under test: an analytic solution or a limiting case, a tiny case
    worked out by hand or with computer algebra, a certified reference
    value, or a published value.
  \Catches Wrong formulas, constants, signs, and conventions, on the inputs
    where the answer is known.
  \Misses Everything away from those inputs. Known answers exist mostly for
    small or simple inputs, the easy corner of the input space. A reference
    that uses another convention, such as a variance divided by $n$ rather
    than $n-1$, raises false alarms.
  \Cost Low when the answer is in a table, higher when it has to be derived.
  \PairsWith Hand-picked examples (\ref{sec:examples}) and edge cases
    (\ref{sec:edges}); real data (\ref{sec:real-data}) with certified
    results; every scope. Generated inputs (\ref{sec:generated}) only where
    a formula gives the answer for every input.
  \Example The sample variance of 1, 2, 3, and 4, which is $5/3$ with the
    divisor $n-1$ and $5/4$ with the divisor $n$.
  \See \TestingRef{Know the answer before you ask}; \DesignRef{No magic
    numbers}.
\end{typecard}



\subsection{Invariant}
\label{sec:invariant}

\begin{typecard}
  \Definition A property that every correct result has, whatever the input,
    such as a conservation law, a normalization, a bound, a symmetry, or a
    consistency relation.
  \Catches Errors that break the structure of a result: probabilities that
    do not sum to one, a negative variance, mass that is lost, a distance
    matrix that is not symmetric. It needs no known answer, so it can be
    checked on any input.
  \Misses A wrong result that still has the property: probabilities that sum
    to one can still be the wrong probabilities.
  \Cost Low once the property is known, and fast to check.
  \PairsWith Generated inputs (\ref{sec:generated}), as a property-based
    test; simulated (\ref{sec:simulated}) and real data
    (\ref{sec:real-data}); end-to-end tests (\ref{sec:end-to-end}), as counts
    and totals that every step preserves.
  \Example In a simulation of a closed system, the total energy stays
    constant, within the error of the integrator, over every time step.
  \See \TestingRef{Test invariants, not only examples}; \DesignRef{Wear belt
    and braces}, on checking the same invariants on every run.
\end{typecard}



\subsection{Metamorphic relation}
\label{sec:metamorphic}

\begin{typecard}
  \Definition A relation between the results for two related inputs: the
    input is changed in a way whose effect the method determines, and the
    result changes accordingly. Examples are reordering, rescaling,
    duplicating, and splitting the data, and monotonicity in a parameter.
  \Catches Errors that leave each single result plausible, such as swapped
    axes, identifiers misaligned after a sort, a unit converted twice, or a
    result that depends on the order of the rows.
  \Misses Errors that respect the relation: a constant bias survives every
    reordering.
  \Cost Moderate: finding relations takes knowledge of the method, and the
    code runs twice. A relation that holds only approximately, for example
    because reordering changes the rounding of a sum, needs a tolerance of
    its own.
  \PairsWith Generated inputs (\ref{sec:generated}), simulated data
    (\ref{sec:simulated}), and real data (\ref{sec:real-data}), because it
    needs no known answer; every scope.
  \Example The coordinates of a set of points, converted from meters to
    kilometers: their estimated spread falls by a factor of 1000, and
    permuting the points leaves it unchanged.
  \See \TestingRef{Change the input, predict the output}.
\end{typecard}



\subsection{Round trip}
\label{sec:round-trip}

\begin{typecard}
  \Definition An operation followed by its inverse returns the original
    input: a file written and read back, a transform and its inverse, a
    unit converted and converted back, a value encoded and decoded. It is a
    special metamorphic relation (\ref{sec:metamorphic}).
  \Catches Information lost on the way, such as precision, units,
    identifiers, or missing values; wrong encodings; and axes swapped in one
    of the two directions.
  \Misses Errors that the two directions share and that cancel: a transform
    with the wrong normalization returns its input perfectly if its inverse
    has the matching error.
  \Cost Low where both directions exist, and fast.
  \PairsWith Generated inputs (\ref{sec:generated}) and edge cases
    (\ref{sec:edges}), such as empty tables and missing values; integration
    tests (\ref{sec:integration}) across a file format; real data
    (\ref{sec:real-data}) in that format.
  \Example A table with missing values, units, and sample identifiers,
    written to the package's file format, read back, and compared with the
    original.
  \See \DesignRef{Nothing lost in translation}.
\end{typecard}



\subsection{Second opinion}
\label{sec:second-opinion}

\begin{typecard}
  \Definition The result is compared with that of an independent
    implementation of the same method: a naive one written for the test,
    such as a plain loop instead of vectorized code, a dense matrix instead
    of a sparse one, or a direct sum instead of a fast transform; or an
    established package, possibly in another language.
  \Catches Errors of optimized code, such as wrong indices in vectorized or
    sparse code and approximations of fast algorithms that go wrong, on every
    input where the two implementations disagree.
  \Misses Errors that both implementations share, because they come from the
    same derivation, the same paper, or the same person. A disagreement shows
    that one of the two is wrong, not which.
  \Cost A naive implementation is quick to write but slow, so it works on
    small inputs only. An established package adds a test dependency, and a
    new release of it can change its numbers.
  \PairsWith Generated inputs (\ref{sec:generated}), because it provides an
    expected value for every input; hand-picked examples
    (\ref{sec:examples}); unit tests (\ref{sec:unit}).
  \Example The product of two small random sparse matrices, compared with
    the product of the same matrices stored as dense ones.
  \See \TestingRef{Get a second opinion}.
\end{typecard}



\subsection{Parameter recovery}
\label{sec:recovery}

\begin{typecard}
  \Definition Data are simulated from the model with known parameters, the
    method is applied to them, and the test checks that it recovers the
    parameters with honest uncertainty: over many simulated datasets, the
    estimates scatter as much as the reported standard errors say, and the
    intervals contain the true values at their nominal rate, their coverage.
  \Catches Bias, standard errors and intervals that are too narrow, tests
    that reject too often under the null hypothesis, and failures in hard
    cases, such as small samples, parameters at the boundary of their range,
    and weak signals.
  \Misses Errors that the simulator shares with the method, and failures on
    data that depart from the model.
  \Cost High: an independent simulator, bounds derived from the sampling
    error, and many replicates, which make studies of coverage slow.
  \PairsWith Simulated data (\ref{sec:simulated}), which it needs, or real
    data (\ref{sec:real-data}) with a planted signal; integration
    (\ref{sec:integration}) and end-to-end tests (\ref{sec:end-to-end}).
  \Example A thousand datasets simulated from a mixed model with known
    variance components: the 95\% intervals contain the true values in about
    95\% of them, within the sampling error of a thousand replicates.
  \See \TestingRef{Plant the answer, then find it}; \TestingRef{Flaky tests
    cry wolf}, on bounds for statistical claims.
\end{typecard}



\subsection{Ground truth}
\label{sec:ground-truth}

\begin{typecard}
  \Definition The result is compared with a truth about real data that was
    established without the method: by the design of an experiment, by an
    independent measurement, or by expert annotation. Examples are truth
    sets for variant calling, spike-in controls of known concentration, mock
    communities of known composition, imaging phantoms, and annotated images.
    The comparison is usually a summary, such as a sensitivity or an error
    rate within a bound. A certified result of a computation is a known
    answer (\ref{sec:known-answer}) instead: it says what the computation
    gives, not what is true.
  \Catches Failures that only real data show, such as those caused by noise,
    artifacts, and departures from the model, and changes in performance
    after a change of the code.
  \Misses Errors smaller than the uncertainty of the truth, and errors where
    the truth is unknown; truth sets often cover only the regions where it is
    certain. A failure may be a limit of the method rather than a bug, since
    a correct implementation can miss the truth: the test checks the
    implementation when its bound comes from the known performance of the
    method, while whether the method suits the data is validation.
  \Cost High: data with a known truth are rare, expensive to make, and often
    large, which makes slow tests, and their license has to allow
    redistribution.
  \PairsWith Real data (\ref{sec:real-data}), which it needs; end-to-end
    tests (\ref{sec:end-to-end}).
  \Example A variant caller run on a region of a reference genome with a
    published truth set, with its sensitivity and precision compared with
    the values published for the method.
  \See \TestingRef{About this guide}, on testing and validation;
    \TestingRef{Know the answer before you ask}, on tolerances derived from
    the precision of a reference.
\end{typecard}



\subsection{Order of convergence}
\label{sec:convergence}

\begin{typecard}
  \Definition A discretization is run at a sequence of resolutions against
    an exact solution, and the rate at which its error falls is compared
    with the order of the method. Where no exact solution exists, the method
    of manufactured solutions provides one. The same idea applies to Monte
    Carlo estimates, whose error falls as $N^{-1/2}$ in the number of
    samples, and to iterative solvers.
  \Catches Bugs that lower the order of the method, such as a boundary
    condition applied one step late or a wrong coefficient, which still give
    a small error on a fine enough grid and so pass a test at one
    resolution.
  \Misses Errors in terms that the exact solution does not exercise, such as
    a term with a derivative that vanishes for the chosen solution, and
    errors of the model itself: the right equations, solved correctly, can
    still be the wrong physics.
  \Cost High: several runs, some at fine resolution, at resolutions fine
    enough for the leading term of the error to dominate and coarse enough
    for rounding error not to.
  \PairsWith Unit (\ref{sec:unit}) and integration tests
    (\ref{sec:integration}) of solvers, quadrature rules, and
    finite-difference schemes, on a problem with an exact or a manufactured
    solution.
  \Example Halving the step of a second-order solver for an ordinary
    differential equation with a known solution divides the error by about
    four.
  \See \TestingRef{Check the order of convergence}.
\end{typecard}



% ==========================================================
\section{Tests by name}
\label{sec:by-name}
% ==========================================================

Many tests are known by a name that is not a type of its own. Some names
stand for a combination of types, such as property-based test; others add
the purpose a test serves or the moment it runs, such as regression test and
smoke test. Table~\ref{tab:by-name} lists common ones, with the types that
they usually combine and what the name adds. The numbers refer to the
subsections of this document.

Some names belong to activities that are not tests of results. Mutation
testing and code coverage check the test suite itself (\TestingRef{A test
that cannot fail tests nothing}), benchmarks measure performance
(\DesignRef{Size matters}), and validation asks whether the method fits the
scientific question, which no test can answer (\TestingRef{About this
guide}).

\begingroup
\small
\begin{xltabular}{\textwidth}{@{}P{0.15}P{0.13}P{0.15}P{0.16}L@{}}
  \caption{Common names of tests, with the types that they usually combine.}
  \label{tab:by-name} \\
  \toprule
  Name & Scope & Inputs & Oracle & What the name adds \\
  \midrule
  \endfirsthead
  \toprule
  Name & Scope & Inputs & Oracle & What the name adds \\
  \midrule
  \endhead
  \bottomrule
  \endlastfoot
  Benchmark problem
    & End-to-end (\ref{sec:end-to-end})
    & A standard problem of the field (\ref{sec:examples})
    & Known answer (\ref{sec:known-answer}), or the results of other codes
      (\ref{sec:second-opinion})
    & A problem that a community agreed on, such as a shock tube in fluid
      dynamics, so that codes can be compared. Not a benchmark of
      performance. \\
  \addlinespace
  Doc test
    & Unit (\ref{sec:unit})
    & The example in the documentation (\ref{sec:examples})
    & Its printed output, as a snapshot (\ref{sec:snapshot}), or it runs
      (\ref{sec:it-runs})
    & Keeps the examples of the documentation working: doctest in Python
      compares their printed output, and R CMD check runs the examples of an
      R package. \\
  \addlinespace
  Fuzz test
    & Unit or end-to-end
    & Generated (\ref{sec:generated}), often malformed
    & It runs (\ref{sec:it-runs}), or the documented error
    & Searches for inputs that crash the code or slip past its validation;
      runs for hours, apart from the test suite. \\
  \addlinespace
  Golden master, characterization test, approval test
    & End-to-end (\ref{sec:end-to-end})
    & Hand-picked examples (\ref{sec:examples}) or real data
      (\ref{sec:real-data})
    & Snapshot (\ref{sec:snapshot})
    & Records the behavior of existing code, often before code without tests
      is changed. In an approval test, a person approves each new reference
      output. \\
  \addlinespace
  Gradient check
    & Unit (\ref{sec:unit})
    & Hand-picked or generated points (\ref{sec:generated})
    & Second opinion (\ref{sec:second-opinion}): finite differences
    & Checks a derivative, coded by hand or computed by automatic
      differentiation, against finite differences of the function. \\
  \addlinespace
  Manufactured solution
    & Unit or integration (\ref{sec:integration})
    & A problem built from a chosen solution
    & Order of convergence (\ref{sec:convergence})
    & Provides an exact solution where none is known: the chosen function,
      inserted into the equations, gives the source terms that it
      requires. \\
  \addlinespace
  Property-based test
    & Unit (\ref{sec:unit})
    & Generated (\ref{sec:generated})
    & Invariant (\ref{sec:invariant}), metamorphic relation
      (\ref{sec:metamorphic}), or round trip (\ref{sec:round-trip})
    & A framework, such as Hypothesis for Python or quickcheck for R,
      generates the inputs, searches the edges, and reduces a failing input
      to a minimal one. \\
  \addlinespace
  Regression test
    & Any
    & The input that exposed a bug
    & Any; often a known answer (\ref{sec:known-answer})
    & Added when a bug is fixed, so that the bug cannot return unnoticed.
      Some communities also call snapshot tests regression tests. \\
  \addlinespace
  Simulation study
    & Integration or end-to-end
    & Simulated data (\ref{sec:simulated}), many datasets
    & Parameter recovery (\ref{sec:recovery})
    & Measures bias and coverage over many simulated datasets; slow, so it
      runs apart from the fast tests. Its coverage is that of intervals, not
      of code. \\
  \addlinespace
  Smoke test
    & End-to-end (\ref{sec:end-to-end})
    & A small hand-picked example (\ref{sec:examples})
    & It runs (\ref{sec:it-runs})
    & Runs first, after a build, an installation, or a change of
      environment, to decide whether the other tests are worth running. \\
\end{xltabular}
\endgroup



\end{document}
